
Consequently, if \((M)\) denotes the family of faithfully flat
morphisms locally of finite presentation, then \(R\) is
\((M)\)-effective. By \SGARefLink{sga3:iv:corollary:6:3:3}{Expos\'e IV, 6.3.3}, \((Y,p)\) therefore
represents the quotient sheaf of \(X_0\) by \(R\) for the fppf
topology, and the base-change assertions follow from \SGARefLink{sga3:iv:proposition:3:4:3:1}{Expos\'e IV,
3.4.3.1}.

\medskip
\phantomsection\label{rem:sga3-V-5.1}%
\noindent\textbf{Remark 5.1.}---%
\footnote{\label{note:sga3-V-30}Editor's note: this paragraph has been added.}
Keep the hypotheses and notation of \hyperref[thm:sga3-V-4.1]{Theorem~4.1}, and suppose in
addition that \(S\) is locally Noetherian and that
\(\pi_0:X_0\to S\) is quasi-projective. We shall prove that
\(\pi:Y\to S\) is quasi-projective.

The preceding hypotheses imply that \(Y\to S\) is of finite type; see
the proof of \hyperref[lem:sga3-V-6.1-ii]{6.1(ii)}. Let \(\mathscr A\) be an invertible
\(\mathcal O_{X_0}\)-module which is ample for \(\pi_0\). By EGA~II,
6.1.12, \(p_\ast\mathscr A\) is an invertible
\(p_\ast\mathcal O_{X_0}\)-module. Hence \(Y\) has an affine open
cover \((V_i)_{i\in I}\) such that \(\mathscr A\) is trivial over each
of the saturated affine open subsets
\[
  U_i=p^{-1}(V_i).
\]

For each \(i\), put
\[
  A_{i,0}=\mathcal O_{X_0}(U_i),
\]
and let \(A_{i,1}\) be the ring of the affine open subset
\(d_0^{-1}(U_i)=d_1^{-1}(U_i)\) of \(X_1\). Let
\(\delta_{i,0}\) (respectively \(\delta_{i,1}\)) be the homomorphism
\(A_{i,0}\to A_{i,1}\) induced by \(d_0\) (respectively \(d_1\)), and
put
\[
  B_i=\mathcal O_Y(V_i)
  =\{b\in A_{i,0}\mid\delta_{i,0}(b)=\delta_{i,1}(b)\}.
\]

Following EGA~II, \S6.5, consider the invertible
\(\mathcal O_{X_0}\)-module
\[
  N_{d_1}\!\left(d_0^\ast\mathscr A\right),
\]
the norm, relative to the finite locally free morphism
\(d_1:X_1\to X_0\), of the invertible
\(\mathcal O_{X_1}\)-module \(d_0^\ast\mathscr A\). Suppose that
\(\mathscr A\), relative to the cover \((U_i)_{i\in I}\), is given by
transition functions
\[
  c_{ij}\in\mathcal O_{X_0}(U_i\cap U_j)^\times.
\]
Then \(N_{d_1}(d_0^\ast\mathscr A)\) is given by the transition
functions
\[
  N_{\delta_1}(\delta_0(c_{ij}))
  \in\mathcal O_{X_0}(U_i\cap U_j)^\times.
\]
By \hyperref[proof:sga3-V-4.1-a]{4(a)}, these elements belong to
\(\mathcal O_Y(V_i\cap V_j)^\times\), and hence define an invertible
\(\mathcal O_Y\)-module \(\mathscr L\) such that
\[
  p^\ast\mathscr L=N_{d_1}(d_0^\ast\mathscr A).
\]
Moreover, for every \(n\in\mathbb N^\ast\),
\[
  p^\ast(\mathscr L^n)
  =N_{d_1}\!\left(d_0^\ast(\mathscr A^n)\right);
\]
see \emph{loc. cit.}, (6.5.2.1).

We now show that \(\mathscr L\) is ample for \(\pi:Y\to S\).
Replacing \(S\) by an affine open subset, we may suppose that \(S\)
is affine. Let \(y\in Y\), choose \(x\in X_0\) with \(p(x)=y\), let
\(V\) be an affine open neighborhood of \(y\) in \(Y\), and put
\(U=p^{-1}(V)\). Since \(\mathscr A\) is \(\pi_0\)-ample, there exist
\(n\in\mathbb N^\ast\) and a section
\[
  s\in\Gamma(X_0,\mathscr A^n)
\]
such that
\[
  x\in(X_0)_s\subseteq U.
\]
With the notation above, \(s\) is represented by sections
\[
  a_i\in A_{i,0}=\mathcal O_{X_0}(U_i),
  \qquad a_i=c_{ij}a_j\ \text{ on }\ U_i\cap U_j,
\]
and \((X_0)_s\) is the union of the open subsets
\[
  U'_i=\{\mathfrak p\in\operatorname{Spec}(A_{i,0})
    \mid a_i\notin\mathfrak p\}.
\]

For every \(i\), put
\[
  N(a_i)=N_{\delta_1}(\delta_0(a_i))\in B_i.
\]
\hyperref[thm:sga3-V-4.1-i]{Theorem~4.1(i)} and
Lemma~\hyperref[lem:sga3-V-4.1.1]{4.1.1} give
\[
  p(U'_i)
  =p\,d_1d_0^{-1}(U'_i)
  =p\,d_1\!\left(
    \{\mathfrak q\in\operatorname{Spec}(A_{i,1})
      \mid\delta_{i,0}(a_i)\notin\mathfrak q\}
    \right),
\]
while
\[
  d_1\!\left(
    \{\mathfrak q\in\operatorname{Spec}(A_{i,1})
      \mid\delta_{i,0}(a_i)\notin\mathfrak q\}
    \right)
  =
  \{\mathfrak p\in\operatorname{Spec}(A_{i,0})
    \mid N_{\delta_1}(\delta_{i,0}(a_i))\notin\mathfrak p\}.
\]
It follows that
\[
  p(U'_i)
  =
  \{\mathfrak p\in\operatorname{Spec}(B_i)
    \mid N(a_i)\notin\mathfrak p\}.
\]
Consequently,
\[
  p((X_0)_s)=Y_{N(s)},
\]
where \(N(s)\) is the section of \(\mathscr L^n\) on \(Y\) defined by
the sections \(N(a_i)\in\mathcal O_Y(V_i)\). Thus
\begin{equation}
  y\in p((X_0)_s)=Y_{N(s)}\subseteq p(U)=V.
  \tag{*}\label{eq:sga3-V-5.1-star}
\end{equation}
This proves that \(\mathscr L\) is ample for \(\pi:Y\to S\), and
hence completes the proof that \(\pi\) is quasi-projective.
